\documentclass[10pt,a4paper]{amsart}
\usepackage{amsmath,amssymb,amsthm,graphicx}
% Retained minimal preamble for the Stacks Project chapter sites-cohomology.tex (Cohomology on Sites).
% Verbatim excerpt of the chapter preamble: Xy-pic setup (native preamble lines 14-18) and the native
% theorem declarations with the mathematical macros (native preamble lines 154-217). Copied mechanically;
% nothing added, removed or redefined. The wrapper class and the amsmath/amssymb/amsthm/graphicx base
% packages are supplied by the assembler, so the class/externaldocument part of the native preamble is omitted.
% For commutative diagrams we use Xy-pic
\usepackage[all]{xy}

% We use 2cell for 2-commutative diagrams.
\xyoption{2cell}
% Theorem environments.
%
\theoremstyle{plain}
\newtheorem{theorem}[subsection]{Theorem}
\newtheorem{proposition}[subsection]{Proposition}
\newtheorem{lemma}[subsection]{Lemma}

\theoremstyle{definition}
\newtheorem{definition}[subsection]{Definition}
\newtheorem{example}[subsection]{Example}
\newtheorem{exercise}[subsection]{Exercise}
\newtheorem{situation}[subsection]{Situation}

\theoremstyle{remark}
\newtheorem{remark}[subsection]{Remark}
\newtheorem{remarks}[subsection]{Remarks}

\numberwithin{equation}{subsection}

% Macros
%
\def\lim{\mathop{\mathrm{lim}}\nolimits}
\def\colim{\mathop{\mathrm{colim}}\nolimits}
\def\Spec{\mathop{\mathrm{Spec}}}
\def\Hom{\mathop{\mathrm{Hom}}\nolimits}
\def\Ext{\mathop{\mathrm{Ext}}\nolimits}
\def\SheafHom{\mathop{\mathcal{H}\!\mathit{om}}\nolimits}
\def\SheafExt{\mathop{\mathcal{E}\!\mathit{xt}}\nolimits}
\def\Sch{\mathit{Sch}}
\def\Mor{\mathop{\mathrm{Mor}}\nolimits}
\def\Ob{\mathop{\mathrm{Ob}}\nolimits}
\def\Sh{\mathop{\mathit{Sh}}\nolimits}
\def\NL{\mathop{N\!L}\nolimits}
\def\CH{\mathop{\mathrm{CH}}\nolimits}
\def\proetale{{pro\text{-}\acute{e}tale}}
\def\etale{{\acute{e}tale}}
\def\QCoh{\mathit{QCoh}}
\def\Ker{\mathop{\mathrm{Ker}}}
\def\Im{\mathop{\mathrm{Im}}}
\def\Coker{\mathop{\mathrm{Coker}}}
\def\Coim{\mathop{\mathrm{Coim}}}

% Boxtimes
%
\DeclareMathSymbol{\boxtimes}{\mathbin}{AMSa}{"02}

%
% Macros for moduli stacks/spaces
%
\def\QCohstack{\mathcal{QC}\!\mathit{oh}}
\def\Cohstack{\mathcal{C}\!\mathit{oh}}
\def\Spacesstack{\mathcal{S}\!\mathit{paces}}
\def\Quotfunctor{\mathrm{Quot}}
\def\Hilbfunctor{\mathrm{Hilb}}
\def\Curvesstack{\mathcal{C}\!\mathit{urves}}
\def\Polarizedstack{\mathcal{P}\!\mathit{olarized}}
\def\Complexesstack{\mathcal{C}\!\mathit{omplexes}}
% \Pic is the operator that assigns to X its picard group, usage \Pic(X)
% \Picardstack_{X/B} denotes the Picard stack of X over B
% \Picardfunctor_{X/B} denotes the Picard functor of X over B
\def\Pic{\mathop{\mathrm{Pic}}\nolimits}
\def\Picardstack{\mathcal{P}\!\mathit{ic}}
\def\Picardfunctor{\mathrm{Pic}}
\def\Deformationcategory{\mathcal{D}\!\mathit{ef}}
\usepackage{cleveref}
\crefname{originallemma}{original lemma}{original lemmas}
\Crefname{originallemma}{Original lemma}{Original lemmas}
\makeatletter
\newlabel{lemma-RHom-into-K-injective}{{0A94}{source}{}{sourcepointer-1}{}}
\newlabel{lemma-RHom-into-K-injective@cref}{{[originallemma][0][]0A94}{[1][0][]source}{}{}{}}
\makeatother
\title{The internal hom complex is well defined up to quasi-isomorphism in both variables}
\author{The Stacks Project Authors}
\date{stacks-project:sites-cohomology.tex:local-snapshot-20261011; GNU FDL 1.2 or later}
\begin{document}
\maketitle
\noindent\textit{Editorial note.} All delivered mathematical text is the author's own, copied from the pinned Stacks Project chapter file sites-cohomology.tex; no statement or proof step was rewritten, reordered or removed.\par\medskip
\noindent\textit{Editorial note.} The proof is the author's complete argument: both H\textasciicircum{}0 groups are computed and shown equal, which is exactly the asserted quasi-isomorphism after the usual reading of the displayed statement.\par\medskip
\noindent\textit{Editorial note.} Cross-chapter restatement: the same label occurs in chapter cohomology for a ringed space. That statement is the special case of this ringed-site statement, so per the lane-s4 rule only this more general occurrence is counted here.\par\medskip
\noindent\textit{Editorial note.} The single labelled reference occurring inside the counted spans is carried in assembly.external\_references as a non-load-bearing navigation pointer to the same chapter.\par\medskip
\noindent\textit{Editorial note.} The preamble is the chapter's own retained minimal preamble (sources/stacks-sites-cohomology/retained-preamble.tex), a verbatim excerpt of the native preamble: the Xy-pic setup and the native theorem declarations with the mathematical macros.\par\medskip
\noindent\textit{Allowed standard prerequisites (editorial).}\begin{itemize}
\item SheafHom complex of complexes of O-modules and its cohomology sheaves. Same chapter, the construction of SheafHom\textasciicircum{}bullet for complexes of modules on a ringed site. Supplies the two complexes compared in the statement and the H\textasciicircum{}0 groups computed in the proof.
\item K-injective complex of O-modules. Same chapter, the definition of K-injectivity and its use in computing derived homs. Supplies the hypothesis that both I\textasciicircum{}bullet and (I')\textasciicircum{}bullet are K-injective, which is what makes the H\textasciicircum{}0 groups computable.
\item Sheaf associated to the presheaf U -> Hom\_\{D(O\_U)\}(L|\_U, M|\_U). Same chapter, the description of the sheaf of homomorphisms in the derived category by restriction to O\_U. Supplies the common value of the two H\textasciicircum{}0 groups, from which quasi-isomorphism follows.
\end{itemize}\medskip
\ifdefined\hypertarget\hypertarget{sourcepointer-1}{}\fi
\noindent\textit{Source reference.} Original lemma 0A94: sources/stacks-sites-cohomology/pinned-original-source.tex, source lines 8720--8720. Same-chapter navigation pointer invoked at proof line 8774, where the author identifies both H\textasciicircum{}0 groups with the same sheafified presheaf of derived homs. The pointer only resolves that reference for the reader; the cited result is not reproduced inside the counted span, and no bibliography citation is involved.\par\medskip
\begin{lemma}
\label{lemma-RHom-well-defined}
Let $(\mathcal{C}, \mathcal{O})$ be a ringed site. Let
$(\mathcal{I}')^\bullet \to \mathcal{I}^\bullet$
be a quasi-isomorphism of K-injective complexes of $\mathcal{O}$-modules.
Let $(\mathcal{L}')^\bullet \to \mathcal{L}^\bullet$
be a quasi-isomorphism of complexes of $\mathcal{O}$-modules.
Then
$$
\SheafHom^\bullet(\mathcal{L}^\bullet, (\mathcal{I}')^\bullet)
\longrightarrow
\SheafHom^\bullet((\mathcal{L}')^\bullet, \mathcal{I}^\bullet)
$$
is a quasi-isomorphism.
\end{lemma}

\begin{proof}
Let $M$ be the object of $D(\mathcal{O})$ represented by
$\mathcal{I}^\bullet$ and $(\mathcal{I}')^\bullet$.
Let $L$ be the object of $D(\mathcal{O})$ represented by
$\mathcal{L}^\bullet$ and $(\mathcal{L}')^\bullet$.
By Lemma \ref{lemma-RHom-into-K-injective}
we see that the sheaves
$$
H^0(\SheafHom^\bullet(\mathcal{L}^\bullet, (\mathcal{I}')^\bullet))
\quad\text{and}\quad
H^0(\SheafHom^\bullet((\mathcal{L}')^\bullet, \mathcal{I}^\bullet))
$$
are both equal to the sheaf associated to the presheaf
$$
U \longmapsto \Hom_{D(\mathcal{O}_U)}(L|_U, M|_U)
$$
Thus the map is a quasi-isomorphism.
\end{proof}

\end{document}
